\section{Introduction}

In the Czech Republic, the planned \Gls{DGR} in crystalline rock will benefit from a geologically stable, predominantly low-permeability host formation that provides an important natural barrier. However, excavation of access tunnels and disposal wells will inevitably alter the surrounding rock mass. Irreversible damage may occur in the form of newly created microcracks caused by blasting and excavation-induced stress redistribution in the immediate vicinity of the excavation; this zone is commonly referred to as the \gls{EDZ}. In addition, a much larger region---the \gls{EIZ}---may experience reversible changes in hydraulic properties. Both the \gls{EDZ} and \gls{EIZ} may act as preferential pathways for contaminant transport to major geological faults.

The methodology proposed by \cite{BrezinaMethodology2023} was developed to characterize \gls{EIZ}-related changes using pore-pressure monitoring during and after excavation. Bayesian inference was subsequently applied to estimate permeability values for intact rock and the \gls{EIZ}, including associated uncertainties. The near-field contaminant transport model was then used to predict concentration at a nearby geological fault serving as a safety indicator and to assess several bounding scenarios of \gls{EIZ} extent and source distance in a preliminary uncertainty analysis.

This study extends that methodology using updated \gls{DGR} design parameters and recent experimental data from \URFB. The present article focuses on the transport sensitivity analysis, while detailed excavation modelling, bentonite review material, and pore-pressure experiment interpretation are retained in the repository as supporting material for the ongoing paper restructuring.

The main objectives are:
\begin{itemize}
    \item Predict the evolution of the rock-mass stress state and key transport parameters, in particular permeability, within the \gls{EIZ} after repository closure and full saturation of the tunnel backfill.
    \item Evaluate radionuclide transport through disposal wells and backfilled access tunnels, accounting for transport within the \gls{EIZ} and the surrounding fracture network.
    \item Quantify uncertainty in the concentration-based safety indicator with respect to transport-model inputs, including fracture-network, \gls{EIZ}, intact-rock, and bentonite parameters.
\end{itemize}
